\cotefile{157-3}
\scanpage{162}
\mspage{111}
% archive sheet number: 162
\uline{Corollaire 2} (Moyennant \unsure{$W(7)$|$W\,7)$}) Soit $i:U\unsure{-|\to}X$ une imm. ouverte $W$-\unsure{dewyérienne|dwyérienne}. Pour que $i$ soit une $W$-cofibration triviale, il f. et s. que $i\in W$, ou ce qui revient au $=$, (\deleted{\unsure{prop.}} cf page 82) que $\dot T\hookrightarrow T$ soit dans $W$.

C'est évidemment nécessaire, et la suffisance résulte du th.~3 a).

\uline{Questions ouvertes} a) J'ignore si une immersion ouverte $i$ qui est une $W$-cofibration, voire une $W$-cofibration triviale, est forcément $W$-dewyérienne.\textsuperscript{*}
% slanted
\margin{\textsuperscript{*} et peut-être $\unsure{=}$ une imm. \unsure{Dwyer}!} Mais par \unsure{le|les} cor. de la prop 7 (p.~85), je sais que c'est le cas si on suppose $T$ isomorphe à la catégorie finale $e$. Il s'ensuit \deleted{\unsure{donc}} au moins ceci : Pour que non seulement $i:U-X$, mais aussi \deleted{\unsure{toute}} les immersions ouvertes $\unsure{i_1|u_1}:U-X_1$ qui s'en déduisent par le procédé \unsure{du|de} prop 7 b) au moyen de \deleted{\unsure{foncteurs}}\inserted{flèches} quelconques $T_1-T$ de but $T$ dans $\mathrm{Cat}$, soient des $W$-cofibrations, il f. et s. que \unsure{$i$|$u$} soit dewyérienne.

b) : \unsure{Y a-t-il} une bonne notion d'\uline{immersion \unsure{dwyérienne|dewyérienne}}, applicable à \uline{toutes} les immersions
